

%%%%%%%%%%%%%%%
\begin{frame}{Une application de messagerie}


On collecte les \alert{besoins} des utilisateurs 

\begin{redbullet}
  \item ``Je veux pouvoir envoyer un message à n'importe quel utilisateur'' 
   \item ``Je veux pouvoir avoir plusieurs destinataires''
  \item ``Je veux savoir qui a envoyé, qui a reçu, quel message''
  \item ``Je veux pouvoir répondre à un message en le citant''
  \item ``Je veux pouvoir retrouver tous les messages suivant des critères comme le 
    nom du destinataire, ou un mot du sujet''
\end{redbullet}


\vfill
Et encore: je veux envoyer des fichiers, masquer l'expéditeur, etc., etc.

\end{frame}

%%%%%%%%%%%%%%%
\begin{frame}{Conception: premières entités}

Rôle du concepteur\,: définir les entités et leurs associations.

\vfill

\figSlideWithSize{ea-messagerie-1.pdf}{8cm}


\vfill
\begin{blocgauche}
Il faut \alert{nommer} les entités, définir leur \alert{identifiant}
et les \alert{cardinalités} des associations.
\end{blocgauche}
\end{frame}


%%%%%%%%%%%%%%%
\begin{frame}{Conception: ajout d'association(s)}

On complète le schéma avec l'avancement de la conception.

\vfill

\figSlideWithSize{ea-messagerie-2.pdf}{8cm}


\vfill
Associations plusieurs-plusieurs\,: on pourrait 
les qualifier. Exemple\,: mode d'envoi (principal, copie
copie cachée, ...)
\end{frame}

%%%%%%%%%%%%%%%
\begin{frame}{Conception: finalisation (?) }

Ici, on ajoute la notion de successeur/prédecesseur aux messages.

\vfill

\figSlideWithSize{ea-messagerie.pdf}{9.5cm}
\vfill


\alert{On va s'en tenir là.} Il faut, si possible, anticiper
les évolutions futures.
\end{frame}


%%%%%%%%%%%%%%%
\begin{frame}{Imaginer une instance}

Pour vérifier, on peut toujours dessiner une instance du modèle.

\vfill

\figSlideWithSize{instance-messagerie.pdf}{9cm}

\vfill


\alert{Encore mieux}\,: On réalise une maquette de la future application 


\end{frame}


\begin{frame}{À retenir}

Conception d'une base = une affaire d'expérience. Constats\,:

\vfill

\begin{redbullet}
\item Faire valider la conception par les utilisateurs régulièrement
  \item Faire simple\,: associations binaires, identifiants séquentiels
  \item Comprendre ce qu'on va pouvoir / ne pas pouvoir représenter
  \item Faire des maquettes, des graphes, des dessins pour vérifier qu'on se comprend bien
\end{redbullet}

\vfill

Délicat  car les besoins ne sont pas toujours très clairs.
\end{frame}

